\documentclass[twoside,leqno]{article}
\usepackage[letterpaper]{geometry}
\usepackage{siamproceedings}
\usepackage[utf8]{inputenc}
\usepackage[T1]{fontenc}
\usepackage[english]{babel}

%\usepackage{authblk}
\usepackage{tabularx}
\usepackage{mdframed}
% SIAM supplies the default 10pt font, amsmath, and ntheorem.
\DeclareMathOperator{\symdiff}{\triangle}
\usepackage{amssymb}
\usepackage{stmaryrd}
\usepackage{mathrsfs} 
\usepackage{mathtools}
\usepackage{thm-restate}
% hyperref and cleveref are configured by siamproceedings.
\usepackage{url}
\usepackage{todonotes}
\usepackage{mathrsfs} 
\usepackage{soul}
\usepackage{relsize}
%\usepackage{enumerate}
\usepackage{enumitem}
\usepackage[all,defaultlines=3]{nowidow}
\usepackage[mathlines]{lineno}
\usepackage{multicol}
\usepackage{enumitem}
\usepackage{xspace}
\usepackage{lineno}
\usepackage{nicefrac}
\usepackage{wasysym}
\usepackage{comment}
%\usepackage{bbold}
%\linenumbers
\usepackage[algo2e,linesnumbered,ruled]{algorithm2e}
% Preserve algorithm2e despite SIAM loading the algorithm package.
\expandafter\let\expandafter\algorithm\csname algorithm2e\endcsname
\expandafter\let\expandafter\endalgorithm\csname endalgorithm2e\endcsname
\usepackage{textpos}
\setlist[itemize]{topsep=4pt,itemsep=3pt,parsep=0pt} 
\setlist[enumerate]{topsep=4pt,itemsep=3pt,parsep=0pt} 
 %\linenumbers


\usetikzlibrary{decorations.markings}

\crefname{algocf}{Algorithm}{Algorithms}
\crefname{claim}{Claim}{Claims}
\crefname{figure}{Figure}{Figures}
\renewcommand{\preceq}{\preccurlyeq}
\renewcommand{\succeq}{\succcurlyeq}
\renewcommand{\subset}{\subseteq}
% SIAM supplies theorem, lemma, corollary, proposition and definition.
% Preserve the manuscript's shared section counter and independent counters.
\renewtheorem*{theorem*}{Theorem}
\newtheorem{conjecture}{Conjecture}
\newtheorem{remark}{Remark}
\newsiamthm{observation}{Observation}
\newtheorem{case}{Case}
\newsiamthm{claim}{Claim}
\newsiamthm{fact}{Fact}
\newtheorem{property}{Property}
\theorembodyfont{\normalfont}
\renewtheorem{definition}[theorem]{Definition}
\newtheorem{example}{Example}
\renewtheorem*{example*}{Example}
\theorembodyfont{\normalfont\itshape}

% Global figure/table numbers, as required by the official SIAM example.
\counterwithout{figure}{section}
\counterwithout{table}{section}
% ntheorem places proof marks automatically, including terminal displays.
\qedsymbol{\proofbox}
\newcommand{\qedhere}{\qed}

\numberwithin{equation}{section}

\renewcommand{\cal}{\mathcal}

\newenvironment{claimproof}[1][Proof of the claim]{%
  \renewcommand{\proofSymbol}{$\lrcorner$}%
  \begin{proof}[#1]%
}{%
  \end{proof}%
}

\DeclareRobustCommand{\stirling}{\genfrac\{\}{0pt}{}}

\newcommand{\drds}[1]{\textsc{Distance-}$#1$ \textsc{Dominating Set}}
\newcommand{\dris}[1]{\textsc{Distance-}$#1$ \textsc{Independent Set}}
\newcommand{\notDominated}[1]{\mathtt{witnessNotDom}(#1)}
\newcommand{\candidate}[1]{\mathtt{candidateDom}(#1)}
\newcommand{\witnessInd}[1]{\mathtt{witnessNotInd}(#1)}
\newcommand{\candidateInd}[1]{\mathtt{candidateInd}(#1)}
\newcommand{\add}[1]{\mathtt{insert}(#1)}
\newcommand{\remove}[1]{\mathtt{remove}(#1)}
\newcommand{\near}[1]{\mathtt{nearVertex}(#1)}
\newcommand{\far}[1]{\mathtt{farVertex}(#1)}
\newcommand{\init}[1]{\mathtt{init}(#1)}
\newcommand{\farVertexDS}[2]{\mathsf{FarVertex}_{#2}[#1]}
\newcommand{\farVertexDSS}[3]{\mathsf{FarVertex}_{#2}^{#3}[#1]}
\newcommand{\nearVertexDS}[2]{\mathsf{NearVertex}_{#2}[#1]}
\newcommand{\nearVertexDSS}[3]{\mathsf{NearVertex}_{#2}^{#3}[#1]}
\newcommand{\dominatingSetDS}[2]{\mathsf{DominatingSet}_{#2}[#1]}
\newcommand{\independentSetDS}[2]{\mathsf{IndependentSet}_{#2}[#1]}
\newcommand{\ds}[1]{\mathtt{domSet}(#1)}
\newcommand{\is}[1]{\mathtt{indSet}(#1)}
\newcommand{\FO}{\mathsf{FO}}
\newcommand{\poly}{\mathrm{poly}}

\newcommand{\avgdeg}{\mathsf{avgdeg}}

\newcommand{\Minors}{\mathsf{Minors}}
\newcommand{\Fra}{\mathsf{Fraternal}}
\newcommand{\fra}{\mathsf{fraternal}}


\newcommand{\ar}{\mathsf{ar}}
\newcommand{\Af}{\mathbb{A}}
\newcommand{\Bf}{\mathbb{B}}

\newcommand{\cc}{\mathsf{cc}}
\newcommand{\bd}{\partial}

\newcommand{\bn}{\mathsf{bn}}
\newcommand{\tw}{\mathsf{tw}}
\newcommand{\wtw}{\mathsf{wtw}}
\newcommand{\bw}{\mathsf{bw}}
\newcommand{\wbw}{\mathsf{wbw}}

\newcommand{\wei}{\mathbf{w}}

\newcommand{\Pp}{\mathcal{P}}
\newcommand{\Qq}{\mathcal{Q}}
\newcommand{\Ww}{\mathcal{W}}

\newcommand{\tup}[1]{{\bar{#1}}}

\newcommand{\Prob}{\mathbb{P}}
\renewcommand{\Pr}{\Prob}
\newcommand{\Exp}{\mathbb{E}}


\newcommand{\clos}{\mathrm{cl}}
\newcommand{\bnd}{\partial}

%\newcommand{\tup}{\bar}
\newcommand{\rt}[2]{#1\langle #2\rangle}


\def\phi{\varphi}
\newcommand{\N}{\mathbb{N}}
\newcommand{\R}{\mathbb{R}}
\newcommand{\Z}{\mathbb{Z}}
\newcommand{\Zp}{\Z^+}
\newcommand{\wh}{\widehat}

\newcommand{\Leaves}{\mathrm{Leaves}}
\newcommand{\Ball}{\mathrm{Ball}}
\newcommand{\dist}{\mathrm{dist}}
\newcommand{\cw}{\mathsf{cvw}}
\newcommand{\wcw}{\mathsf{wcvw}}
\newcommand{\ct}{\mathsf{cut}}
\newcommand{\md}{\mathsf{mid}}
\newcommand{\sli}{\mathsf{sli}}

\newcommand{\Aa}{\mathcal{A}}
\newcommand{\Bb}{\mathcal{B}}
\newcommand{\Ee}{\mathcal{E}}

\newcommand{\Tt}{\mathcal{T}}
\newcommand{\Ll}{\mathcal{L}}
\newcommand{\Ff}{\mathcal{F}}
\newcommand{\Ss}{\mathcal{S}}

\newcommand{\F}{\mathcal{F}}
\newcommand{\G}{\Cc}
\newcommand{\val}{\mathrm{val}}
\newcommand{\col}{\mathrm{col}}

\newcommand{\Hom}{\mathrm{Hom}}
\newcommand{\Sub}{\mathrm{Sub}}
\newcommand{\ISub}{\mathrm{ISub}}


\newcommand{\from}{\colon}
\newcommand{\Cc}{\mathscr C}
\newcommand{\Dd}{\mathscr D}
\def\epsilon{\varepsilon}
\def\eps{\varepsilon}
\newcommand{\Oh}{\mathcal{O}}

\newcommand{\id}{\mathrm{id}}
\renewcommand{\emptyset}{\varnothing}

\newcommand{\michal}[1]{\todo[color=pink!40]{Micha\l{}: #1}}
\newcommand{\michalin}[1]{\todo[size=\normalsize,inline,color=pink!40]{Micha\l{}: #1}}

\newcommand{\wn}[1]{\todo[color=green!501]{{\textbf{W:} #1}}}
\newcommand{\wninline}[2][inline]{\todo[color=green!50,#1]{{\textbf{W:} #2}}}
\newcommand{\bbinline}[2][inline]{\todo[color=blue!50,#1]{{\textbf{B:} #2}}}
\newcommand{\azinline}[2][inline]{\todo[color=yellow!50,#1]{{\textbf{A:} #2}}}
\newcommand{\mpinline}[2][inline]{\todo[color=pink!40,#1]{{\textbf{M:} #2}}}

\newcommand{\WReach}{\textrm{WReach}}
\newcommand{\wcol}{\textrm{wcol}}

\newcommand{\UnMaps}{\mathsf{UnweightedMappings}}
\newcommand{\WeiMaps}{\mathsf{WeightedMappings}}
\newcommand{\MapEx}{\mathsf{MappingExample}}
\newcommand{\RetVer}{\mathsf{RetrieveVertex}}
\newcommand{\WeiSum}{\mathsf{WeightedSum}}
\newcommand{\Ver}{\mathsf{Verify}}
\newcommand{\MapVer}{\mathsf{MapVertex}}


\renewcommand{\leq}{\leqslant}
\renewcommand{\geq}{\geqslant}
\renewcommand{\le}{\leq}
\renewcommand{\ge}{\geq}

\renewcommand{\setminus}{-}
\newcommand{\set}[1]{\{#1\}}
\newcommand{\setof}[2]{\{#1 \colon#2\}}

\newcommand{\cnst}[3]{{#1}_{#2}^{\ref{#3}}}

\newcommand{\aff}[1]{\normalsize\textcolor{black!60}{#1}}

\usepackage{tikz}

\SetKwBlock{Repeat}{repeat}{}
\let\oldnl\nl% Store \nl in \oldnl
\newcommand{\nonl}{\renewcommand{\nl}{\let\nl\oldnl}}

\newcommand{\affi}[1]{\textcolor{black!50}{#1}}

\tikzstyle{vertex}=[circle,draw=black,fill=gray,minimum size=0.15cm,inner sep=0pt]

\tikzstyle{svertex}=[circle,draw=black,fill=black,minimum size=0.11cm,inner sep=0pt]

% Keep only the first URL printed by siamplain (DOI takes precedence).
% Restrict the change to bibliography formatting; retain the BibTeX data.
\makeatletter
\newcommand{\SODAfirstbiburl}[1]{%
  \SODAoriginalurl{#1}%
  \@ifnextchar,{\SODAcheckbiburlcomma}{}%
}
\def\SODAcheckbiburlcomma,{%
  \@ifnextchar\url{\SODAskipbiburl}{,\space}%
}
\def\SODAskipbiburl\url#1{%
  \@ifnextchar,{\SODAcheckbiburlcomma}{}%
}
\AddToHook{env/thebibliography/begin}{%
  \let\SODAoriginalurl\url
  \let\url\SODAfirstbiburl
}
\makeatother

\begin{document}

\newcommand{\funding}{
	\begin{minipage}[t]{0.30\textwidth}
		\includegraphics[scale=0.15]{LOGO_ERC-FLAG_EU.png}
	\end{minipage}\hfill
	\begin{minipage}[b]{0.65\textwidth}\footnotesize
		All the authors were supported by the project BOBR that received funding from the European Research Council (ERC) under the European Union’s Horizon 2020 research and innovation programme, grant agreement No. 948057. In particular, during the work on this manuscript Mi. Pilipczuk was affiliated with the University of Warsaw, and the majority of work of B. Bosek on this manuscript was done while he was also affiliated with the University of Warsaw.
	\end{minipage}}



\title{\Large Dynamic domination and independence in sparse graphs\thanks{\funding}}
%\date{\today}
\date{}
\author{
	Bart\l{}omiej Bosek\thanks{Jagiellonian University, Kraków (\email{bartlomiej.bosek@uj.edu.pl}).}
	\and
	Wojciech Nadara\thanks{University of Warsaw (\email{w.nadara@mimuw.edu.pl}).}
	\and
	Michał Pilipczuk\thanks{Technical University of Darmstadt (\email{pilipczuk@mathematik.tu-darmstadt.de}).}
	\and
	Anna Zych-Pawlewicz\thanks{University of Warsaw (\email{anka@mimuw.edu.pl}).}
}
%\author[1]{Anonymous authors}
%	\author[1,2]{Bart\l{}omiej Bosek}
%	\author[1]{Wojciech Nadara}
%	\author[1]{Micha\l{} Pilipczuk}
%	\author[1]{Anna Zych-Pawlewicz}
%	\affil[1]{\aff{University of Warsaw, Poland}}
%	\affil[2]{\aff{Jagiellonian University, Cracow, Poland}}
\maketitle

% Copyright mode: 0 = unresolved, 1 = SIAM, 2 = authors, 3 = organization.
% Set this only after the authors select the corresponding copyright agreement.
\newcommand{\SODAcopyrightmode}{2}
\fancyfoot[R]{\scriptsize
  \ifcase\SODAcopyrightmode
    SODA 2027: copyright selection pending
  \or
    Copyright \textcopyright\ 2027 by SIAM\\
    Unauthorized reproduction of this article is prohibited
  \or
    Copyright \textcopyright\ 2027\\
    Copyright for this paper is retained by authors
  \or
    Copyright \textcopyright\ 2027\\
    Copyright retained by principal author's organization
  \else
    \PackageError{SODA}{Invalid copyright mode}{Use 0, 1, 2, or 3.}
  \fi
}

\begin{abstract}
	Let $\Cc$ be a class of graphs of bounded expansion and $r,k\in \N$ be fixed. We give a dynamic data structure that for a given dynamic graph $G$, updated by edge insertions and deletions subject to the promise that $G\in \Cc$ at all times, maintains the answer to the following two queries:
	\begin{itemize}[nosep]
		\item Does $G$ contain a distance-$r$ dominating set of size $k$?
		\item Does $G$ contain a distance-$r$ independent set of size $k$?
	\end{itemize}
	The data structure is randomized with error probability bounded by $\eps$, for a parameter $\eps>0$ fixed upon the initialization. The amortized update time is $\log^c n\cdot \log \frac{1}{\eps}$, where $n$ is the vertex count of $G$ and $c$ is a constant that depends only on $r$, $k$, and $\Cc$. In the case of the first query, the data structure can also output a distance-$r$ dominating set of size $k$, if existent.
	
	We also prove that when $r=1$, our data structure for the dominating set query can be implemented even if we only assume that the maintained graph $G$ has degeneracy bounded by a constant $d$, yielding a simpler data structure with an improved amortized update time of $2^{k^{\Oh(d)}}\cdot \log^3 n\cdot \log \tfrac{1}{\eps}$. Finally, we prove that in graphs of degeneracy at most $d$, one can maintain an $\Oh(d^2)$-approximation of the minimum size of a (distance-$1$) dominating set with amortized expected update time $d^{\Oh(1)}\cdot \log n$.
\end{abstract}

% \begin{textblock}{20}(-1.75, 3.6)
% \includegraphics[width=40px]{logo-erc}%
% \end{textblock}
% \begin{textblock}{20}(-1.75, 4.6)
% \includegraphics[width=40px]{logo-eu}%
% \end{textblock}

% Continuous proceedings layout: no separate title page or contents page.

\section{Introduction}\label{sec:intro}

\emph{Sparsity} is a research area in structural graph theory that studies the structure in graphs that exclude dense local obstructions, formalized through the notion of shallow minors. More precisely, we say that a graph $H$ is a \emph{depth-$r$ minor} of a graph $G$ if a supergraph of $H$ can be obtained from $G$ by contracting mutually disjoint connected subgraphs of radius at most $r$. The two main concepts are the following:
\begin{itemize}
	\item A graph class $\Cc$ has \emph{bounded expansion} if for every $r\in \N$ there is a finite upper bound $c_r\in \N$ on the average degree of depth-$r$ minors of graphs from $\Cc$.
	\item A graph class $\Cc$ is \emph{nowhere dense} if for every $r\in \N$ there is a finite upper bound $t_r\in \N$ on the sizes of complete graphs that can be found as depth-$r$ minors of graphs from $\Cc$.
\end{itemize}
Clearly, if $\Cc$ has bounded expansion then $\Cc$ is nowhere dense as well, but the reverse implication does not hold in general. Many well-studied classes of sparse graphs, such as graphs with bounded maximum degree, graphs excluding a fixed (topological) minor (in particular planar graphs), or some natural classes of sparse geometric intersection graphs, do have bounded expansion. Therefore, any methods designed for bounded expansion classes will apply to them as well. We refer the reader to the monograph of Ne\v{s}et\v{r}il and Ossona de Mendez~\cite{sparsity} or to the lecture notes of Pilipczuk, Pilipczuk, and Siebertz~\cite{sparsityNotes} for a broad introduction to the vast toolbox of techniques for classes of bounded expansion and nowhere dense classes.

The methods of Sparsity are particularly useful in the context of parameterized algorithms, as witnessed by the award of the Nerode Prize 2025 to Ne\v{s}et\v{r}il and Ossona de Mendez for laying foundations of the field. Not surprisingly, these methods apply predominantly to parameterized problems of local nature, such as the following:
\begin{itemize}
	\item \textsc{$H$-(Induced)-Subgraph Isomorphism}: Decide whether the fixed graph $H$ is an (induced) subgraph of the given graph $G$.
	\item \drds{r}: Decide whether the given graph $G$ has a \emph{distance-$r$ dominating set} of size $k$, that is, a set $D$ of $k$ vertices such that every vertex of $G$ is at distance at most $r$ from some vertex of $D$.
	\item \dris{r}: Decide whether the given graph $G$ has a \emph{distance-$r$ independent set} of size $k$, that is, a set $I$ of $k$ vertices that are pairwise at distance more than $r$ from each other.
\end{itemize}
All these problems are $\mathsf{W}[1]$-hard on general graphs, but become fixed-parameter tractable on any nowhere dense class of graphs. This is a manifestation of a more general phenomenon. Namely, each of the problems above can be expressed by a sentence of first-order logic $\FO$, of length depending on the relevant parameter(s): the vertex count of $H$, or $r$ and $k$. As proved by Grohe, Kreutzer, and Siebertz~\cite{GroheKS17}, the \emph{model-checking problem} for~$\FO$ --- given a graph $G$ and an $\FO$ sentence $\varphi$, decide whether $\varphi$ holds in $G$ ---  can be solved in almost linear fixed-parameter time on any nowhere dense class $\Cc$; precisely in time $\Oh_{\Cc,\varphi,\eps}(|G|^{1+\eps})$, for any fixed~$\eps>0$, where $|G|$ denotes the number of vertices of $G$. As proved earlier by Dvo\v{r}\'ak, Kr\'al', and Thomas~\cite{DvorakKT13}, on bounded expansion classes one can even obtain a linear fixed-parameter time of $\Oh_{\Cc,\varphi}(|G|)$. While fixed-parameter tractability of \textsc{$H$-(Induced)-Subgraph Isomorphism}, \drds{r}, \dris{r} on nowhere dense classes follows from these meta-theorems, dedicated algorithms for those problems were known already earlier~\cite{DawarK09,NesetrilM08a}. In fact, their development was vital in establishing the toolbox needed~for~\cite{DvorakKT13,GroheKS17}.

Since the fundamental work of Dvo\v{r}\'ak et al.~\cite{DvorakKT13} and of Grohe et al.~\cite{GroheKS17}, the approach to model-checking on classes of sparse graphs has been thoroughly understood, particularly for classes of bounded expansion. The developed framework is robust and has been successfully extended to the settings of enumeration, counting, and other aggregate queries~\cite{GroheS18,KazanaS19,SchweikardtSV22,Torunczyk20}, circuit complexity~\cite{PilipczukST18}, approximation~\cite{Dvorak22}, and distributed computing~\cite{BlinFFGGMRT26}. One setting that has so far resisted progress, despite efforts, is that of dynamic data structures. 

Precisely, consider the following setting. Fix a class of bounded expansion $\Cc$ and an $\FO$ sentence $\varphi$. Let $G$ be a fully dynamic graph that is updated over time by edge insertions and edge deletions (the vertex set stays fixed), subject to a guarantee that $G\in \Cc$ at all times. The goal is to design a data structure that efficiently maintains whether $\varphi$ is satisfied in $G$.

The question about the existence of such a data structure was asked in 2013 by Dvo\v{r}\'ak and T\r{u}ma~\cite{DvorakT13}, who showed that this is indeed possible for the \textsc{$H$-Induced-Subgraph Isomorphism} problem. Precisely, they gave a dynamic data structure that is able to maintain the number of induced copies of $H$ in $G$ (as well as the number of $H$-subgraphs in $G$ and the number of homomorphisms from $H$ to~$G$) with amortized update time $\Oh_{\Cc,h}(\log^{\binom{h}{2}-1} n)$, where $h$ and $n$ denote the vertex counts of $H$ and $G$, respectively. Somewhat curiously, their approach is based on the Inclusion-Exclusion Principle, so even though their data structure can count the number of solutions, it is not capable of actually providing an example solution.

\textsc{$H$-Induced-Subgraph Isomorphism} is a specific example of an $\FO$-definable problem, because the sentence expressing it is purely existential: After quantifying the vertices of $H$   existentially, the adjacencies and nonadjacencies can be verified using a quantifier-free subformula. Consequently, algorithms for this problem on classes of sparse graphs are typically much simpler; see e.g.~\cite{NesetrilM08a}. Since the work of Dvo\v{r}\'ak and T\r{u}ma in 2013, no progress on their question has been reported, even for the more complicated $\exists\forall$ problems such as \drds{r} and \dris{r}.

We remark that some previous works~\cite{DvorakKT13,Torunczyk20} considered the partially dynamic setting where (essentially) the graph stays fixed and its vertices and edges gets recolored, but this setting is much simpler than the fully dynamic one that we study here.

\paragraph*{Our contribution.} In this work we propose fully dynamic data structures for the \drds{r} and \dris{r} problems on classes of bounded expansion, with polylogarithmic amortized update time.\footnote{For a tuple of parameters $\tup p$, the $\Oh_{\tup p}(\cdot)$ notation hides multiplicative factors that may depend on $\tup p$.} 

\begin{restatable}{theorem}{mainDomBE}\label{thm:main-domset-be}
	Fix a graph class $\Cc$ of bounded expansion, $r,k\in \N$, and $\eps>0$. Then there is a randomized dynamic data structure that given a dynamic $n$-vertex graph $G$, updated by edge insertions and removals subject to guarantee that it always belongs to $\Cc$, supports the following query:
	\begin{itemize}
		\item $\ds{}$: return a distance-$r$ dominating set of size at most $k$ in $G$, or $\bot$ if no such set exists.
	\end{itemize}
	Every answer to the query is correct with probability at least $1-\eps$ against an oblivious adversary.
	The amortized time complexity of an update or a query is $\log^{\Oh_{\Cc,r,k}(1)} n\cdot \log \tfrac{1}{\eps}$. The data structure can be initialized on an edgeless graph in time $\Oh_{\Cc,r,k}(n\log n\log \tfrac{1}{\eps})$ and occupies  $\Oh_{\Cc,r,k}(n\log n\log \tfrac{1}{\eps})$ space at all times.
\end{restatable}

\vspace{-0.3cm}

\begin{restatable}{theorem}{mainIndBE}\label{thm:main-indset-be}
	Fix a graph class $\Cc$ of bounded expansion, $r,k\in \N$, and $\eps>0$. Then there is a randomized dynamic data structure that given a dynamic $n$-vertex graph $G$, updated by edge insertions and removals subject to guarantee that it always belongs to $\Cc$, supports the following query:
	\begin{itemize}
		\item $\is{}$: decide whether in $G$ there exists a distance-$r$ independent set of size $k$.
	\end{itemize}
	Every answer to the query is correct with probability at least $1-\eps$ against an oblivious adversary.
	The amortized time complexity of an update or a query is $\log^{\Oh_{\Cc,r,k}(1)} n\cdot \log \tfrac{1}{\eps}$. The data structure can be initialized on an edgeless graph in time $\Oh_{\Cc,r,k}(n\log n\log \tfrac{1}{\eps})$ and occupies  $\Oh_{\Cc,r,k}(n\log n\log \tfrac{1}{\eps})$ space at all times.
\end{restatable}

Note that the data structure of \cref{thm:main-domset-be} is capable of providing an example solution, while that of \cref{thm:main-indset-be} only reports the existence of a solution. Both data structures are randomized against an oblivious adversary, which means that the guarantees on the probability of correctness hold only assuming the adversary cannot choose the next update/query based on the data structure's answers to the previous queries. We stress that the data structure's answers to the consecutive queries are \emph{not} independent random variables; they all depend on some common random bits fixed upon the initialization. 

At this point, the proofs of \cref{thm:main-domset-be,thm:main-indset-be} do not extend to $\FO$-definable problems beyond \drds{r} and \dris{r}. This is because our approach is based on dynamizing two dedicated algorithms for these problems due to Fabia\'nski, Pilipczuk, Siebertz, and Toru\'nczyk~\cite{FabianskiPST19}, obtained via their framework of \emph{progressive exploration}. We believe that this might indicate that in order to resolve the question of Dvo\v{r}\'ak and T\r{u}ma, one has to design a new $\FO$-model-checking algorithm on classes of sparse graphs that would be more in the spirit of progressive exploration.

We remark that on the way to \cref{thm:main-domset-be,thm:main-indset-be}, we resolve the question posed by Dvo\v{r}\'ak and T\r{u}ma~\cite{DvorakT13} by extending their data structure so that it can also report examples of relevant mappings. See \cref{thm:dvorak-tuma-examples} for a precise formulation.

\medskip

Next, we observe that for the specific case $r=1$ (so the classic \textsc{Dominating Set} problem), the data structure of \cref{thm:main-domset-be} can be simplified and works already under the assumption of bounded degeneracy (equivalently, of bounded arboricity or of bounded maximum average degree among subgraphs).

\begin{restatable}{theorem}{mainDomDeg}\label{thm:main-domset-deg}
	Fix $d,k\in \N$ and $\eps>0$. Then there is a randomized dynamic data structure that given a dynamic $n$-vertex graph $G$, updated by edge insertions and removals subject to guarantee that $G$ is $d$-degenerate at all times, supports the following query:
	\begin{itemize}
		\item $\ds{}$: return a dominating set of size at most $k$ in $G$, or $\bot$ if no such set exists.
	\end{itemize} 
	Every answer to the query is correct with probability at least $1-\eps$ against an oblivious adversary.
	The amortized time complexity of an update or a query is $2^{k^{\Oh(d)}}\cdot \log^3 n \log \tfrac{1}{\eps}$. The data structure can be initialized on an edgeless graph in time $k^{\Oh(d)}\cdot n\log n\log \tfrac{1}{\eps}$ and uses  $k^{\Oh(d)}\cdot n\log n\log \tfrac{1}{\eps}$ space at all times.
\end{restatable}

Finally, it is also known that on bounded expansion classes, both \drds{r} and \dris{r} admit constant-factor approximation algorithms~\cite{Dvorak13}. We prove that for $r=1$, this can be lifted to the dynamic setting, even assuming only bounded~degeneracy.

\begin{restatable}{theorem}{mainApx}\label{thm:apx-dynamic}
For every $d\in \N$ there is a randomized dynamic data structure (against an oblivious adversary) that given a dynamic $n$-vertex graph $G$, updated by edge insertions and removals, subject to guarantee that $G$ is $d$-degenerate at all times, maintains a dominating set $D$ in $G$ satisfying $|D|\leq (4d+1)^2\cdot \mathsf{dom}_1(G)$, where $\mathsf{dom}_1(G)$ denotes the minimum size of a dominating set in $G$. The access to $D$ is provided by queries that report its size in time $\Oh(1)$ and enumerate all its vertices in time $\Oh(|D|)$. The data structure can be initialized on an edgeless graph in time $\Oh(n)$, uses $\Oh(dn)$ space at all times, and processes every edge insertion and deletion in expected amortized time $d^{\Oh(1)}\cdot \log n$.
\end{restatable}

Whether \cref{thm:apx-dynamic} can be lifted to $r>1$ under the assumption that the maintained graph belongs to a fixed bounded-expansion class $\Cc$ at all times, remains an interesting open question. Note here that while the setting $r=1$ is trivial for the \dris{r} problem (an $n$-vertex $d$-degenerate graph contains an independent set of size $\frac{n}{d+1}$), it becomes again interesting for $r>1$. In fact, it is known~\cite{Dvorak13} that in bounded-expansion classes, there is a constant-factor multiplicative gap between the minimum size of a distance-$r$ dominating set and the maximum size of a distance-$(2r+1)$ independent set. So if only an approximate value of the optimum needs to be reported, the questions of dynamic approximation for \drds{r} and \dris{(2r+1)} in bounded expansion classes are~equivalent.

\medskip

\paragraph*{Acknowledgements.} We thank \L{}ukasz Kowalik, Piotr Sankowski, and Marek Soko\l{}owski for many discussions over several years that, perhaps indirectly, led to this work.

\medskip

\paragraph*{AI disclosure.} No AI tools were used during the work on this manuscript.

\section{Overview}\label{sec:overview}

In this section we provide a high-level exposition of the technical ideas leading to our findings, particularly to the main results: \cref{thm:main-domset-be,thm:main-indset-be}. A complete presentation of all our results and of the proofs leading to them can be found in the full version of this paper, available on arxiv~\cite{ourArxiv}.

\subsection{Progressive exploration}

The key idea behind \cref{thm:main-domset-be,thm:main-indset-be} is to develop dynamic counterparts of the \emph{progressive exploration} algorithms for \drds{r} and \dris{r} in sparse graphs, proposed by Fabia\'nski, Pilipczuk, Siebertz, and Toru\'nczyk~\cite{FabianskiPST19}. For the sake of concreteness, let us focus our discussion on the \drds{r} problem.

Suppose that we are given a graph $G$ and positive integers $r$ and $k$. We would like to find a distance-$r$ dominating set of size $k$ in $G$, or conclude that no such set exists. A progressive exploration algorithm for this problem --- called by Fabia\'nski et al. the \emph{semi-ladder algorithm} --- proceeds in rounds, and constructs a sequence of \emph{candidates} $D_1,D_2,D_3,\ldots$, which are $k$-tuples of vertices, and \emph{witnesses} $w_1,w_2,w_3,\ldots$, which are single vertices. With $D_1,\ldots,D_{i-1}$ and $w_1,\ldots,w_{i-1}$ having been constructed in rounds $1,2,\ldots,i-1$, the $i$th round proceeds as follows:
\begin{description}
	\item[Step 1:] Determine whether there exists a $k$-tuple of vertices $D_i$ that distance-$r$ dominates all the witnesses found so far, that is, $\{w_1,\ldots,w_{i-1}\}$. If not, then $G$ does not admit a distance-$r$ dominating set of size $k$ and we can terminate the algorithm.
	\item[Step 2:] Verify whether $D_i$ distance-$r$ dominates the whole graph. If so, then $D_i$ is a valid solution that can be reported. Otherwise, select any vertex $w_i$ that is not distance-$r$ dominated by $D_i$, and proceed with the new candidate $D_i$ and witness $w_i$ to the next round.
\end{description}
It is clear that when this algorithm reports an outcome, then this outcome is correct. However, a priori it is not at all clear that the algorithm will reach an outcome within a small number of rounds. This seems particularly problematic, because the task  in Step~1 --- finding a $k$-tuple of vertices that distance-$r$ dominates all the $i-1$ witnesses gathered so far --- gets more and more computationally expensive with every round. (We have not discussed its implementation so far, we will do it shortly.)

Somewhat surprisingly, Fabia\'nski et al. proved that on sparse graphs, more precisely on any nowhere dense class of graphs, the semi-ladder algorithm always terminates within a bounded number of rounds.

\begin{theorem}[follows from~\cite{FabianskiPST19}]\label{thm:overview-semiladder}
	For every $r\in \N$ and every nowhere dense class of graphs $\Cc$, there exists a constant $c$ such that the semi-ladder algorithm deployed on any graph $G\in \Cc$ with parameters $r$ and $k$, terminates after performing at most $k^c$ rounds.
\end{theorem}

See also~\cite{FabianskiPST19} for a more precise formulation, and recall that bounded expansion classes are in particular nowhere dense. We remark that the proof of \cref{thm:overview-semiladder} requires deep tools from the theory of Sparsity, in particular the equivalent characterization of nowhere denseness through the notion of \emph{flatness} (also known as \emph{uniform quasi-wideness}).

\cref{thm:overview-semiladder} makes the semi-ladder algorithm perfect for lifting to the dynamic setting. Namely, since the number of rounds executed by the algorithm is bounded by a constant depending only on the class $\Cc$ and the parameters $r$ and $k$, we can just run the semi-ladder algorithm after every update to the graph, provided we are able to maintain data structures that allow efficient implementation of every~round.

Consider then the $i$th round of the algorithm. The task in Step~2 boils down to a single application of the following query in $G$, with $S=D_i$ (so $|S|\leq k$):
\begin{itemize}
	\item $\far{S,r}$: Given a set $S\subseteq V(G)$, find a vertex $v$ such that $\dist(v,s)>r$ for all $s\in S$, or return $\bot$ if no such vertex exists.
\end{itemize}
On the other hand, Step~1 can be implemented using at most $k^i$ queries of the following kind, with $|S|<i$:
%\wninline{To $k^i$ jest trochę konfundujące. Może zmienić jednak na bardziej naturalne $2^i$? Podobnie w lemacie 4.4}
%\michalin{Mh, słusznie :D. Zaimplementujesz?}
\begin{itemize}
	\item $\near{S,r}$: Given a set $S\subseteq V(G)$, find a vertex $v$ such that $\dist(v,s)\leq r$ for all $s\in S$, or return $\bot$ if no such vertex exists.
\end{itemize}
Indeed, we may iterate through all the partitions of the set of witnesses $\{w_1,\ldots,w_{i-1}\}$ into $k$ sets, and for each partition $\Pp$, invoke $\near{S,r}$ on each part $S$ of $\Pp$ to verify whether $S$ can be \mbox{distance-$r$} dominated by a single vertex $v_S$. Note that if for some partition $\Pp$ all these checks go through, then $D_i\coloneqq \{v_S\colon S\in \Pp\}$ is a $k$-vertex set that distance-$r$ dominates $\{w_1,\ldots,w_{i-1}\}$; and otherwise, if the checks for all the partitions fail, then $\{w_1,\ldots,w_{i-1}\}$ cannot be distance-$r$ dominated by $k$ vertices. There are at most $k^{i-1}$ partitions to check, and checking each requires $k$ applications of the $\near{S,r}$~query.

The discussion above combined with \cref{thm:overview-semiladder} amounts to the proof of the following statement.

\begin{proposition}\label{prop:overview-reduction}
	Let $r,k\in \N$ and $\Cc$ be a nowhere dense class of graphs. Then the semi-ladder algorithm on any graph $G\in \Cc$ with parameters $r$ and $k$ can be implemented using $k^{\Oh_{\Cc,r}(1)}$ queries of the form $\near{S,r}$ with $|S|\leq k^{\Oh_{\Cc,r}(1)}$, and $k^{\Oh_{\Cc,r}(1)}$ queries of the form $\far{S,r}$ with $|S|\leq k$.
\end{proposition}

So from now on our task is to design data structures that can efficiently answer queries $\near{S,r}$ and $\far{S,r}$ with sets $S$ of bounded size and $r$ being a fixed constant. As we will see shortly, our data structure for $\near{S,r}$ is a non-trivial extension of the result of Dvo\v{r}\'ak and T\r{u}ma~\cite{DvorakT13}, which in particular answers one of their open questions. On the other hand, the data structure for $\far{S,r}$ is more involved and requires new insights.

We remark that our data structures for $\near{S,r}$ and $\far{S,r}$ queries heavily rely on the assumption that the fixed class $\Cc$, to which the graph $G$ is guaranteed to belong at all times, has bounded expansion, even though \cref{prop:overview-reduction} actually works in the larger generality of nowhere dense classes. In \cref{sec:conclusions}, we comment on the possibility of lifting our results to nowhere dense classes.

\medskip

Finally, let us discuss the case of \dris{r}. For this problem, Fabia\'nski et al. gave a different algorithm, called the \emph{ladder algorithm}, which is somewhat more complicated than the semi-ladder algorithm and whose proof of correctness requires more insight. However, the bottom line is that on any nowhere dense class $\Cc$ and for fixed parameters $r$ and $k$, the ladder algorithm again terminates within $\Oh_{\Cc,r,k}(1)$ rounds and each of those rounds can be implemented using $\Oh_{\Cc,r,k}(1)$ calls to (slight generalizations of) the queries $\near{S,r}$ and $\far{S,r}$ with $|S|\leq \Oh_{\Cc,r,k}(1)$. So designing  data structures implementing these queries efficiently is sufficient to obtain the data structure for the \dris{s} promised in \cref{thm:main-indset-be}. We remark that in the case of a positive outcome, the ladder algorithm actually does not provide a distance-$r$ independent set of size $k$, but only a proof of its existence; this is why this restraint is also present in \cref{thm:main-indset-be}.
%\wninline{``feature'' -> ``restraint''? Feature wydaje mi się trochę pozytywnie nacechowane.}

\subsection{Finding near vertices}

We start by designing a data structure for the $\near{S,r}$ queries, which is the easier part of the argument. For this, we first recall the result of Dvo\v{r}\'ak and T\r{u}ma~\cite{DvorakT13}, which requires some terminology.

For two graphs $H$ and $G$, a \emph{homomorphism} from $H$ to $G$ is a mapping $\varphi\colon V(H)\to V(G)$ such that for any edge $uv$ of $H$, $\varphi(u)\varphi(v)$ is an edge in $G$. A \emph{subgraph isomorphism} is an injective homomorphism, and an \emph{induced subgraph isomorphism} is one where the implication above is an equivalence: $uv$ is an edge in $H$ if and only if $\varphi(u)\varphi(v)$ is an edge in $G$. We denote by $\Hom(H,G),\Sub(H,G),\ISub(H,G)$ the sets of homomorphisms, subgraph isomorphisms, and induced subgraph isomorphisms from $H$ to $G$, respectively.

With this terminology, the result of Dvo\v{r}\'ak and T\r{u}ma reads as follows.

\begin{theorem}[\cite{DvorakT13}]\label{thm:overview-DT}
	Fix a bounded expansion class $\Cc$, graph $H$, and $\Ff\in \{\Hom,\Sub,\ISub\}$. Then there is a data structure that for a dynamic graph $G$ on $n$ vertices, guaranteed to belong to $\Cc$ at all times, maintains the value $|\Ff(H,G)|$ with amortized update time $\Oh_{\Cc,H}(\log^{\binom{h}{2}-1} n)$, where $h$ is the vertex count of $H$. The data structure can be initialized on an edgeless $G$ in $\Oh_{\Cc,H}(n)$ time and uses $\Oh_{\Cc,H}(n)$ space at all~times.	
\end{theorem}

As already remarked in \cite{DvorakT13}, \cref{thm:overview-DT} is quite robust and can be easily extended to the setting of vertex- and edge-colored graphs, where vertices and edges bear colors (that can be updated), and these colors have to be preserved under homomorphisms and (induced) subgraph isomorphisms. Note that by placing unique colors on some vertices, we can extend \cref{thm:overview-DT} to the setting of \emph{rooted mappings}: we can query for the number of homomorphisms/(induced) subgraph isomorphisms that map some fixed tuple $\tup u$ of vertices of $H$ to a given tuple $\tup v$ of vertices of $G$.

Coming back to the implementation of the query $\near{S,r}$, the following claim is straightforward and translates the query to the language of homomorphisms.

\begin{claim}\label{cl:homs}
	Let $r\in \N$, $G$ be a graph, and $S=\{s_1,\ldots,s_\ell\}$ be a set of vertices of $G$. Then the following conditions are equivalent:
	\begin{itemize}
		\item There is a vertex $v$ in $G$ such that $\dist(v,s_i)\leq r$ for all $i\in \{1,\ldots,\ell\}$.
		\item There is a function $\rho\colon \{1,\ldots,\ell\}\to \{0,1,\ldots,r\}$ such that there is a homomorphism $\varphi$ from the graph $H_\rho$ in \cref{fig:homs} to $G$ satisfying $\varphi(z_i)=s_i$ for all $i\in \{1,\ldots,\ell\}$.
	\end{itemize}
\end{claim}

\begin{figure}
	\centering
	\begin{tikzpicture}
		
		\node[vertex] (x) at (5,2) {};
		\node at (5,2.5) {$x$};
		
		\foreach \i/\j in {1/1,2/2,3/3,4/4,5/5,8/\ell-1,9/\ell} {
			\node[vertex] (z\i) at (\i,0) {};
			\node at (\i,-0.5) {$z_{\j}$};
			\draw[
			postaction={decorate, decoration={
					markings,
					mark=at position 0.25 with {\node[svertex] {};},
					mark=at position 0.5 with {\node[svertex] {};},
					mark=at position 0.75 with {\node[svertex] {};}
			}}
			] (x) -- (z\i);
 		}
 	
 		\draw[<->,dashed] (1-0.1,0+0.3) -- (5-0.22,2+0.22); 
		\node at (6.5,0) {$\dots$};
		\node at (3-0.5,1+0.5) {$\rho(1)$};
		
	\end{tikzpicture}
	\caption{Graph $H_\rho$ in \cref{cl:homs}. It is constructed from  vertex $x$ and vertices $z_1,\ldots,z_\ell$ by connecting $x$ with $z_i$ by a path of length $\rho(i)$, for each $i\in \{1,\ldots,\ell\}$.}\label{fig:homs}
\end{figure}

Therefore, in order to obtain a data structure for just the existential query $\near{S,r}$ --- whether a suitable vertex $v$ exists --- it suffices to maintain the data structure of \cref{thm:overview-DT} for every possible graph $H_\rho$ (of which there are $(r+1)^\ell$ many, where $\ell$ is an upper bound on $|S|$), where the vertices $z_1,\ldots,z_\ell$ are marked with unique colors. Upon query $\near{S,r}$, we enumerate the vertices of $S$ as $s_1,\ldots,s_\ell$, mark them using the colors of $z_1,\ldots,z_\ell$, and determine whether any of the maintained data structures indicates that there exists at least one homomorphism from some $H_\rho$ to $G$ that maps each $z_i$ to the respective $s_i$. Once this is determined, vertices $s_i$ can be unmarked.

However, this method only determines the \emph{existence} of a vertex $v$ that is close to all the vertices of $S$, and does not provide an \emph{example} of such a vertex. To implement the semi-ladder algorithm, we actually have to be able to retrieve a suitable vertex $v$, because it should serve as a witness for the next rounds. While we believe that this could be done for homomorphisms by a careful inspection of the arguments in~\cite{DvorakT13}, how to do this for (induced) subgraph isomorphisms is not clear at all, due to an application of the Inclusion--Exclusion Principle; see also the discussion in~\cite[Section~9]{DvorakT13}.

We now present a robust way to mitigate this caveat and turn a counting data structure into an example-reporting data structure with the help of randomization. While for the $\near{S,r}$ query this is probably not necessary, because the approach presented above relies only on homomorphisms, the method that we are going to introduce now will be reused in the implementation of $\far{S,r}$. More precisely, we will use an elegant fingerprint retrieval technique that was recently used in the context of data structures for parameterized problems~\cite{Majewski24,Nadara22}; let us describe it now.

The first observation is that \cref{thm:overview-DT} can be easily lifted to count \emph{weighted} mappings. Precisely, we assume that we additionally have a fixed weight function $\wei\colon V(H)\times V(G)\to \Z$, and the contribution of a mapping $\varphi$ to the count is $\val_\wei(\varphi)\coloneqq \prod_{u\in V(H)} \wei(u,\varphi(u))$, instead of just $1$. Precisely, instead of $|\Ff(H,G)|$, the data structure maintains the value
\[\val_\wei(\Ff(H,G))\coloneqq \sum_{\varphi\in \Ff(H,G)} \val_\wei(\varphi).\]
Suppose now that we have a fixed vertex $x$ of $H$ and we would like to find, if existent, any vertex $v$ of $G$ such that there exists $\varphi\in \Ff(H,G)$ with $v=\varphi(x)$. Suppose further, for a moment, that there is at most one such vertex $v$. By assuming without loss of generality that the vertex set of $G$ consists of numbers $\{1,\ldots,n\}$, we consider two weight functions $\wei,\wei'\colon V(H)\times V(G)\to \Z$ defined as follows:
\begin{align*}
	\wei(y,w) &\coloneqq 1\qquad \textrm{for all }(y,w)\in V(H)\times V(G);\\
	\wei'(y,w) &\coloneqq \begin{cases}w & \textrm{if }y=x,\\ 1 & \textrm{otherwise.}\end{cases}
\end{align*}
The key observation is the following: if indeed the sought vertex $v$ is unique, then it is equal to the ratio
\[
\frac{\val_{\wei'}(\Ff(H,G))}{\val_{\wei}(\Ff(H,G))}.
\]
Note that this ratio can be easily computed by maintaining the (weight-extended) data structures of Dvo\v{r}\'ak and T\r{u}ma for weight functions $\wei$ and $\wei'$. Observe also that if the denominator of this ratio --- $\val_\wei(\Ff(H,G))$ --- is equal to $0$, then this means that there is no $\Ff$-mapping from $H$ to $G$ and consequently no vertex $v$ with the required property.

The remaining question is how to lift the assumption that there is only at most one vertex $v$ that can be the image of $x$ in a mapping $\varphi\in \Ff(H,G)$. We do this with the help of randomization. Let us fix a desired bound $\eps>0$ on the error probability. Upon initialization of the data structure, for each $i\in \{0,1,\ldots,\lceil \log n\rceil\}$ we sample $\xi=\Theta(\log \tfrac{1}{\eps})$ vertex subsets $S_{i,1},\ldots,S_{i,\xi}\subseteq V(G)$, where each set $S_{i,j}$ is constructed by including every vertex independently with probability $\frac{1}{2^{i+1}}$, so that the expected size of $S_{i,j}$ is $\tfrac{n}{2^{i+1}}$. It is not hard to see that if the set $X\coloneqq \{v\mid \textrm{there is }\varphi\in \Ff(H,G)\textrm{ such that }\varphi(x)=v\}$ is non-empty, then with probability at least $1-\eps$ at least one of the sets $S_{i,j}$ will have intersection of size exactly $1$ with $X$. Therefore, we may apply the same trick as in the previous paragraph, but for every set  $S_{i,j}$ simultaneously, where in the data structures constructed for $S_{i,j}$ we nullify all the weights of the pairs $(y,w)$ with $y=x$ and $w\notin S_{i,j}$. Thus, with probability at least $1-\eps$ at least one of the data~structures points to an example vertex $v\in X$, whose feasibility can be verified using the data structure of~\cref{thm:overview-DT}.

The fingerprinting retrieval technique presented above can be applied iteratively to retrieve example images of the vertices of $H$ under a homomorphism or an (induced) subgraph isomorphism one by one, until a full example mapping is retrieved. In this way, we answer the open question of Dvo\v{r}\'ak and T\r{u}ma~\cite{DvorakT13} about extending their data structure with a possibility to report example (induced) copies of a fixed graph $H$ in a dynamic graph $G$ that always belongs to a fixed class of bounded expansion. Precisely, we prove the following result. (A \emph{$k$-colored graph} is one where each edge is colored using one of $k$ colors.)

\begin{restatable}{theorem}{reporting} \label{thm:dvorak-tuma-examples}
	Let $\eps>0$, $H$ be a fixed graph, $\G$ be a graph class of bounded expansion, $k\in \N$, and $\F$ be either $\Hom$, $\Sub$, or $\ISub$. 
	Let $G$ be a dynamic $k$-colored
	graph on $n$ vertices that belongs to $\G$ at all times.
	Then, there is a randomized data structure $\MapEx_{\F, H, k,\eps}[G]$
	which after every update is able to report that either $\F(H, G)$ is empty, or provide some $\phi \in \F(H, G)$. The amortized update time is $\Oh_{\G, H, k}(\log^{\Oh_H(1)} n \log{\frac{1}{\eps}})$, the initialization time is $\Oh_{\G, H, k}(n \log n \log {\tfrac{1}{\eps}})$, and the space complexity is $\Oh_{\G, H, k}(n \log n \log {\tfrac{1}{\eps}})$.  The data structure never provides false positives, but may fail to provide an example mapping $\phi$ with probability at most $\eps$, against an oblivious adversary.
\end{restatable}

\newcommand{\fr}{\mathsf{far}}

\subsection{Finding far vertices}

We now proceed to the description of the data structure supporting the query $\far{S,r}$. For simplicity, we shall focus on a simpler task: We would like to be able to report the quantity $|V_\fr|$, where
\[V_\fr\coloneqq \{v\in V(G)\mid \dist(v,s)>r\textrm{ for all }s\in S\}.\]
Once we achieve this, a data structure that can actually report an example vertex $v\in V_\fr$ can be obtained using the fingerprint retrieval technique described in the previous section, as follows:
\begin{itemize}
	\item First, we generalize the data structure so that it reports the quantity $\sum_{v\in V_\fr} \wei(v)$, for a weight function $\wei\colon V(G)\to \Z$ fixed upon initialization.
	\item Then, we apply the fingerprint retrieval technique so that by maintaining the data structures from the first point for $\Oh(\log n\log \tfrac{1}{\eps})$ different weight functions $\wei$ allows us to retrieve the index of an example vertex from $V_\fr$, provided $V_\fr\neq \emptyset$.
\end{itemize}

Furthermore, we will assume that in all the queries, $S$ consists of a single fixed vertex $s$, $S=\{s\}$. This can be achieved by adding a fresh isolated vertex $s$ to the graph that at the time of a query, is made adjacent to all the vertices of $S$. Then the original query $\far{S,r}$ is equivalent to the query $\far{\{s\},r+1}$ after the modification.

At this point we reach the crucial novel part of this work. Namely, counting $|V_\fr|$ cannot be directly formulated as counting homomorphisms, because roughly speaking, $V_\fr$ is defined by the \emph{non-existence} of a homomorphism. Instead, we will compute $|V_\fr|$ using the Inclusion-Exclusion Principle, by adding and subtracting vertices that are, in various ways, close to $s$. As there may be an unbounded number of short paths connecting any pair of vertices, it is a priori unclear how to apply the Inclusion-Exclusion Principle over them. We overcome this issue by exploiting the  properties of bounded expansion graph classes again.
Precisely, we use the same Sparsity tool as Dvo\v{r}\'ak and T\r{u}ma~\cite{DvorakT13}: \emph{fraternal augmentations}.

% \wninline{Może jakieś zdanie wyjaśnienia czemu to nietrywialne? Np. ``As there may be an unbounded number of short paths connecting any pair of vertices, it is a priori unclear how to apply the Inclusion-Exclusion Principle over them, which is why we again exploit properties of bounded expansion graph classes.''?}
% \michalin{Dobry pomysł, zaimplementowałem i pozmieniałem lekko strukturę}

%To this end, 

\paragraph*{Fraternal augmentations.}
Let $G$ be a graph belonging to a fixed bounded expansion class $\Cc$. The assumption that~$G\in \Cc$ implies that for some constant $d_1=\Oh_\Cc(1)$, $G$ is \emph{$d_1$-degenerate}: every subgraph of $G$ has a vertex of degree at most $d_1$. Hence, by iteratively deleting vertices of the smallest degree from $G$ and orienting incident edges outwards, we obtain an orientation $\vec G_1$ with maximum outdegree at most~$d_1$.

Now, we add \emph{fraternal edges} to $\vec G_1$: for every pair of (directed) edges $(u,v)$ and $(u,v')$ with a common tail $u$, we add an (undirected) edge $vv'$, unless it was already present. It turns out that the new graph $G_2$ --- consisting of the oriented edges of $G$ and of the added fraternal edges --- is still sparse: the underlying undirected graph belongs to some graph class $\Cc_2$, depending only on $\Cc$, that still has bounded expansion. Consequently, the new fraternal edges can be oriented so that we obtain an oriented supergraph $\vec G_2$ of $\vec G_1$ whose maximum outdegree is at most $d_1+d_2$, for some constant $d_2$ depending only on $\Cc_2$. By performing this operation $r$ times, we eventually obtain an orientation $\vec G_r$ of a supergraph of $G$ whose maximum outdegree is at most $d\coloneqq d_1+d_2+\ldots+d_r$. 

With every edge of $\vec G_r$ we may naturally associate its \emph{length}: the original edges of $G$ have length $1$, and a fraternal edge added to edges of lengths $a$ and $b$ is assigned length $a+b$. Edges of length larger than $r$ will have no significance for us, so we may just not create them in the process.

An important idea in the work of Dvo\v{r}\'ak and T\r{u}ma~\cite{DvorakT13} is that the sequence of fraternal augmentations described above can be maintained in the dynamic setting with amortized update time $\Oh_{\Cc,r}(\log^r n)$. The key component here is the classic data structure of Brodal and Fagerberg~\cite{BrodalF99}, which can be used to maintain the orientation at each level of the construction. In this ``tower'' of Brodal--Fagerberg data structures, every update to $G$ triggers $\Oh_{\Cc}(\log n)$ updates to $\vec G_1$ in the amortized sense, which in turn triggers $\Oh_{\Cc}(\log^2 n)$ updates to $\vec G_2$, and so on. 

So from now on we assume that we maintain the augmentation $\vec G_r$ in our data structure, where every edge is decorated with its length. Moreover, since the outdegree of every vertex is at most $d$, we may additionally maintain a labelling of the edges of $\vec G_r$ with labels $\{1,\ldots,d\}$ so that the edges with the same tail are assigned pairwise different labels.


\begin{figure}
	\centering
	\begin{tikzpicture}
		
		\node[vertex] (z0) at (0,0) {};
		\node[vertex] (z15) at (15,0) {};
		\draw (z0) -- (z15);
		
		\foreach \i in {1,2,3,4,5,6,7,8,9,10,11,12,13,14} {
			\node[svertex] (z\i) at (\i,0) {};
		}
		
		\draw[->,very thick, red] (z0) edge[bend left] (z3);
		\draw[->,very thick, red] (z3) edge[bend left] (z5);
		
		\draw[->, very thick, red] (z5) edge (z6);
		\draw[->, very thick, red] (z6) edge[bend left] (z8);
		
		\draw[->, very thick, blue] (z15) edge[bend right] (z13);
		\draw[->, very thick, blue] (z13) edge[bend right] (z10);
		\draw[->, very thick, blue] (z10) edge[bend right] (z8);
		
		\node at (0,-0.5) {$a=w_{\alpha_0}$};
		\node at (3,-0.5) {$w_{\alpha_1}$};
		\node at (5,-0.5) {$w_{\alpha_2}$};
		\node at (6,-0.5) {$w_{\alpha_3}$};
		\node at (8,-0.5) {$w_{\alpha_4}=w_{\beta_3}$};
		\node at (10,-0.5) {$w_{\beta_2}$};
		\node at (13,-0.5) {$w_{\beta_1}$};
		\node at (15,-0.5) {$w_{\beta_0}=b$};
		
	\end{tikzpicture}
	\caption{An $a$-$b$-path with a shortcut. The directed paths $Q,R$ of the shortcut are resp. red and blue.}\label{fig:shortcut}
\end{figure}


\paragraph*{Shortcuts.}
Suppose now that in $G$ we have two vertices $a$ and $b$ that can be connected by a path $P=(a=w_0,w_1,\ldots,w_{r'}=b)$ of some length $r'\leq r$. If we trace what happens with $P$ during the consecutive augmentations, then it is not hard to see that in $\vec G_r$, $P$ will have an oriented ``shortcut'' consisting of:
\begin{itemize}
	\item a directed path $Q$ of the form $w_{\alpha_0}\to w_{\alpha_1}\to \ldots \to w_{\alpha_p}$, where $0=\alpha_0<\alpha_1<\ldots<\alpha_p$, so that each edge $(w_{\alpha_j},w_{\alpha_{j+1}})$ has length at most $\alpha_{j+1}-\alpha_j$; and
	\item a directed path $R$ of the form $w_{\beta_0}\to w_{\beta_1}\to \ldots \to w_{\beta_q}$, where $r'=\beta_0>\beta_1>\ldots>\beta_q=\alpha_p$, so that each edge $(w_{\beta_j},w_{\beta_{j+1}})$ has length at most $\beta_j-\beta_{j+1}$.
\end{itemize}
See \cref{fig:shortcut}. If we now record the lengths and the labels of the consecutive edges of the shortcut paths $Q$ and $R$, we obtain two sequences of pairs from $\{1,\ldots,d\}\times \{1,\ldots,r\}$ such that the sum of all the second coordinates is at most $r$; we call such a pair of sequences the \emph{pattern} of a shortcut. At this point, two observations are crucial:
\begin{enumerate}[label=(O\arabic*)]
	\item\label{o:numpat} There are only at most $(dr)^{2r}=\Oh_{\Cc,r}(1)$ different patterns. Let $\Pi$ be their set.
	\item\label{o:unique} For every pair of vertices $s,u\in V(G)$ and pattern $\pi=(\pi_Q,\pi_R)$, there is at most one shortcut of an $a$-$b$-path with pattern $\pi$. This is because in $\vec G_r$, the edges outgoing from a single vertex have pairwise different labels, so knowing $a$, $b$, and sequences $\pi_Q,\pi_R$, we can uniquely reconstruct the paths $Q$ and $R$. A shortcut is present if the reconstructed paths $Q$ and $R$ end at the same vertex.
\end{enumerate}
Based on the above discussion, the following is now clear.

\begin{claim}\label{cl:overview-shortcuts}
	For any two vertices $a,b\in V(G)$, the following conditions are equivalent:
	\begin{itemize}
		\item $\dist(a,b)\leq r$.
		\item In $\vec G^r$, there exists a shortcut between $a$ and $b$ with some pattern $\pi\in \Pi$.
	\end{itemize}
\end{claim}

We remark that the concept of path shortcuts in the context of fraternal augmentations that we describe above actually dates back to the 2006 work of Kowalik and Kurowski~\cite{KowalikK06}.

\paragraph*{Counting far vertices.} We may now come back to our initial goal of maintaining the cardinality of the set $V_\fr$, which consists of vertices $u$ with $\dist(s,u)>r$. For a vertex $u$ and a pattern $\pi\in \Pi$, call $\pi$ \emph{realized} at $u$ if in $\vec G_r$ there is a shortcut between $s$ and $u$ with pattern $\pi$. By \cref{cl:overview-shortcuts}, $V_\fr$ consists of those vertices $u$ at which no pattern of $\Pi$ is realized. Now, if for $\Gamma\subseteq \Pi$ we define $V_\Gamma$ as the set of those $u\in V(G)$ at which all the patterns of $\Gamma$ are realized, then by the Inclusion-Exclusion principle, we have
\[|V_\fr|=\sum_{\Gamma\subseteq \Pi} (-1)^{|\Gamma|}\cdot |V_\Gamma|.\]
(Note here that $V_\emptyset=V(G)$.) Since $\Pi$ is of size $\Oh_{\Cc,r}(1)$ by \ref{o:numpat}, to compute $|V_\fr|$ it suffices to compute $|V_\Gamma|$ for each $\Gamma\subseteq \Pi$. However, it is easy to argue using \ref{o:unique} that $|V_\Gamma|$ is equal to the number of homomorphisms $\varphi$ from the oriented graph $H_\Gamma$ depicted in \cref{fig:HGamma} to $\vec G_r$ that satisfy $\phi(z)=s$. (Here, we consider both $H_\Gamma$ and $\vec G_r$ as directed graphs with edges decorated with lengths and labels, and the considered homomorphisms have to respect those decorations.) It now remains to note that this homomorphism count can be maintained using the data structure of \cref{thm:overview-DT}. So all in all, we maintain the data structures of \cref{thm:overview-DT} for all graphs $H_\Gamma$ for $\Gamma\subseteq \Pi$, and from the maintained homomorphism counts we piece together the value of $|V_\fr|$ using Inclusion-Exclusion. This concludes the description of the data structure supporting query $\far{S,r}$ in classes of bounded expansion.




\subsection{Other results}\label{sec:overview-other}

Finally, let us comment on the two side results, \cref{thm:main-domset-deg,thm:apx-dynamic}.

\medskip

The proof of \cref{thm:main-domset-deg} follows the same path as that of \cref{thm:main-domset-be}, except that for $r=1$, queries $\near{S,1}$ and $\far{S,1}$ can be implemented in a much simpler way, and relying only on the boundedness of degeneracy. Let us consider the query $\near{S,1}$, which boils down to finding a vertex $v$ that is simultaneously adjacent to all the vertices of $S$. Since the maintained graph $G$ is $d$-degenerate, we may use the data structure of Brodal and Fagerberg~\cite{BrodalF99} to maintain its orientation $\vec G$ with maximum outdegree $\Oh(d)$. Now, when searching for a suitable vertex $v$, we may first test every vertex in the outneighborhood of $S$ in $\vec G$; this amounts to testing $\Oh(d|S|)$ vertices, each in time $\Oh(d|S|)$. If we do not find a suitable $v$ in this way, we know that we are looking for $v$ such that all the edges between $v$ and $S$ are oriented away from $v$ in $\vec G$. The idea now is that together with $\vec G$, we can efficiently maintain the hypergraph of the outneighborhoods of all the vertices (which are of size $\Oh(d)$), together with all their subsets (which amounts to $2^{\Oh(d)}\cdot n$ subsets in total). Then verifying whether there exists a vertex $v$ with $S$ contained in its neighborhood amounts to checking whether $S$ belongs to the maintained hypergraph.

The argument presented above easily extends to reporting a suitable vertex $v$ adjacent to all the vertices of $S$, as well as counting the number of such vertices $v$. Similarly, for every subset $S'\subseteq S$, we may count the number of vertices $v$ that are adjacent to all the vertices of $S'$. Thus, using the Inclusion-Exclusion Principle, we may also count the number of vertices $v$ that are adjacent to none of the vertices in $S$. This can be easily extended to weighted counting (under any weight function $\wei\colon V(G)\to \Z$ fixed upon initialization), so we may again use the fingerprint retrieval technique to design a randomized data structure that can report a vertex non-adjacent to all the vertices of $S$, provided there exists one. This gives an implementation of $\far{S,1}$ and completes the proof of \cref{thm:main-domset-deg}.

\medskip

Finally, the proof of \cref{thm:apx-dynamic} mostly relies on completely different ideas. Similarly as above, using the data structure of Brodal and Fagerberg~\cite{BrodalF99} we may maintain an orientation $\vec G$ of the maintained graph $G$ with maximum outdegree $\Oh(d)$. The key idea is that if we consider the set system of closed outneighborhoods $\Ss\coloneqq \{\{v\}\cup N^+_{\vec G}(v)\colon v\in V(G)\}$, then for any inclusion-wise maximal packing of disjoint sets ${\cal M}\subseteq \Ss$, $\bigcup {\cal M}$ is an $\Oh(d^2)$-approximation of the minimum dominating set. This observation is in essence already present in an old distributed $\Oh(d^2)$-approximation algorithm for \textsc{Dominating Set} on degenerate graphs, due to Lenzen and Wattenhofer~\cite{LenzenW10}. It is not hard to see that the set system $\Ss$ can be maintained efficiently under updates to $\vec G$, hence we can use the recent data structure of Assadi and Solomon~\cite{AssadiSolomon21} to efficiently maintain an inclusion-wise maximal packing in $\Ss$ as well.

In our argumentation we take an extra mile to clarify that the combinatorial argument behind the approximation guarantee is in fact more general, as it extends to larger $r$. Namely, if we assume that $G$ belongs to a fixed class of bounded expansion $\Cc$, then there exists a vertex ordering $\sigma$ of $G$ whose \emph{weak $r$-coloring number} $\wcol_r(G,\sigma)$ is bounded by a constant $c$ depending only on $\Cc$ and $r$. This means that all the \emph{weak $r$-reachability sets} $\WReach_r[G,\sigma,v]$, for $v\in V(G)$, are of size bounded by $c$; these are analogues of the closed outneighborhoods in the $r=1$ case. We observe that again, the union of any inclusion-wise maximal packing in the set system $\Ff_r(G,\sigma)\coloneqq \{\WReach_r[G,\sigma,v]\colon v\in V(G)\}$ is a $c^2$-approximation of the minimum distance-$r$ dominating set in $G$. Unfortunately, for $r>1$ we do not know how to efficiently maintain $\Ff_r(G,\sigma)$ under updates to $G$, as a single update to $G$ may necessitate an unbounded number of changes to $\Ff_r(G,\sigma)$.
\begin{figure}
	\centering
	\begin{tikzpicture}
		
		\node[vertex] (z) at (-4,0) {};
		\node at (-4.5,0) {$z$};
		
		\node[vertex] (x) at (4,0) {};
		\node at (4.5,0) {$x$};
		
		\foreach \i in {1,2,3,4,5} {
			\node[vertex] (u\i) at (0,3-\i*1) {};
		}
		
		\foreach \t in {1,2,3} {
			\node[svertex] (z1\t) at (-\t,2-0.5*\t) {};
		}
		\draw[very thick, red,->] (z) -- (z13);
		\draw[very thick, red,->] (z13) -- (z12);
		\draw[very thick, red,->] (z12) -- (z11);
		\draw[very thick, red,->] (z11) -- (u1);
		
		
		\foreach \t in {1,2,3} {
			\node[svertex] (z4\t) at (-\t,-1+0.25*\t) {};
		}
		\draw[very thick, red,->] (z) -- (z43);
		\draw[very thick, red,->] (z43) -- (z42);
		\draw[very thick, red,->] (z42) -- (z41);
		\draw[very thick, red,->] (z41) -- (u4);
		
		\foreach \t in {2} {
			\node[svertex] (z2\t) at (-\t,1-0.25*\t) {};
		}
		\draw[very thick, red,->] (z) -- (z22);
		\draw[very thick, red,->] (z22) -- (u2);
		
		\foreach \t in {2} {
			\node[svertex] (z5\t) at (-\t,-2+0.5*\t) {};
		}
		\draw[very thick, red,->] (z) -- (z52);
		\draw[very thick, red,->] (z52) -- (u5);
		
		
		\foreach \t in {1,2} {
			\node[svertex] (z3\t) at (-\t*1.33,0) {};
		}
		\draw[very thick, red,->] (z) -- (z32);
		\draw[very thick, red,->] (z32) -- (z31);
		\draw[very thick, red,->] (z31) -- (u3);
		
		
		\foreach \t in {1,2} {
			\node[svertex] (x1\t) at (\t*1.33,2-0.66*\t) {};
		}
		\draw[very thick, blue,->] (x) -- (x12);
		\draw[very thick, blue,->] (x12) -- (x11);
		\draw[very thick, blue,->] (x11) -- (u1);
		
		
		\foreach \t in {1,2} {
			\node[svertex] (x2\t) at (\t*1.33,1-0.33*\t) {};
		}
		\draw[very thick, blue,->] (x) -- (x22);
		\draw[very thick, blue,->] (x22) -- (x21);
		\draw[very thick, blue,->] (x21) -- (u2);
		
		
		\foreach \t in {1,2,3} {
			\node[svertex] (x3\t) at (\t,0) {};
		}
		\draw[very thick, blue,->] (x) -- (x33);
		\draw[very thick, blue,->] (x33) -- (x32);
		\draw[very thick, blue,->] (x32) -- (x31);
		\draw[very thick, blue,->] (x31) -- (u3);
		
		
		\foreach \t in {1,2,3} {
			\node[svertex] (x5\t) at (\t,-2+\t*0.5) {};
		}
		\draw[very thick, blue,->] (x) -- (x53);
		\draw[very thick, blue,->] (x53) -- (x52);
		\draw[very thick, blue,->] (x52) -- (x51);
		\draw[very thick, blue,->] (x51) -- (u5);
		
		\foreach \t in {2} {
			\node[svertex] (x4\t) at (\t,-1+\t*0.25) {};
		}
		\draw[very thick, blue,->] (x) -- (x42);
		\draw[very thick, blue,->] (x42) -- (u4);
		
		
	\end{tikzpicture}
	\caption{Example graph $H_\Gamma$ constructed for a set $\Gamma\subseteq \Pi$ of size $5$. Every pair (red path, blue path) represents a shortcut with some pattern $\pi\in \Gamma$. Note that both the red edges and the blue edges are decorated with $\{1,\ldots,d\}\times \{1,\ldots,r\}$ (not depicted).}\label{fig:HGamma}
\end{figure}
\section{Conclusions}\label{sec:conclusions}

Let us conclude by discussing a handful of open questions.
\begin{itemize}
	\item The central open question is that of Dvo\v{r}\'ak and T\r{u}ma~\cite{DvorakT13}: Can one design a dynamic data structure for $\FO$ model-checking on any fixed  class of bounded expansion with polylogarithmic amortized update time? A natural approach to this question is to try to develop dynamic counterparts for various techniques of Sparsity underlying the known static algorithms~\cite{DvorakKT13,grohe2011methods,PilipczukST18}, such as transitive--fraternal augmentations, weak coloring numbers, or low treedepth colorings. After a few years of attempts, we consider this route hopeless. Instead, we believe that developing a new model-checking algorithm for $\FO$ on classes of bounded expansion, which would be more in the spirit of progressive exploration, could be a way to go.
	\item While the data structure of Dvo\v{r}\'ak and T\r{u}ma~\cite{DvorakT13} can be also deployed on any nowhere dense graph class, with the resulting amortized update time becoming $\Oh_{\Cc,H,\delta}(n^\delta)$ for any $\delta>0$, this is not the case for our data structures of \cref{thm:main-domset-be,thm:main-indset-be}. The reason is that in the nowhere dense setting, the fraternal augmentations discussed in \cref{sec:overview} have maximum outdegree $\Oh_{\Cc,r,\delta}(n^\delta)$ instead of a constant, and hence the Inclusion--Exclusion formula used there may involve $2^{\Oh_{\Cc,r,\delta}(n^\delta)}$ terms, which is superpolynomial. We leave it open whether \drds{r} and \dris{r} on nowhere dense classes admit dynamic data structures with amortized update time $\Oh_{\Cc,r,k,\delta}(n^\delta)$, for any $\delta>0$. However, we remark that the size of this sum is the only reason why this argument does not lift to nowhere dense classes and multiple other results, in particular \cref{thm:dvorak-tuma-examples}, can be lifted verbatim.
	\item Can our data structures be derandomized? At this point, we crucially rely on randomization in the fingerprint retrieval technique, to turn counting data structures into example-reporting data structures.
	\item It remains open whether the result of \cref{thm:apx-dynamic} can be extended to distance-$r$ domination for $r>1$ under the assumption that $G$ belongs at all times to a fixed class of bounded expansion $\Cc$. As discussed in \cref{sec:overview-other}, this is connected to the dynamic maintenance of the set system of weak $r$-reachability sets.
\end{itemize}


%\wninline{Czemu referencja do Dvoraka z WADS wygląda tak biednie?}
\bibliographystyle{siamplain}
\bibliography{references}


\end{document}
